%\documentclass{beamer}
%\usepacakge[UTF8]{ctex}
%----------------------------------------------------------------------------------------
%    PACKAGES AND THEMES
%----------------------------------------------------------------------------------------

\documentclass[aspectratio=169,xcolor=dvipsnames]{beamer}
\usepackage[UTF8]{ctex}
\usetheme{SimplePlus}
\setbeamertemplate{footline}{}

\usepackage{hyperref}
\usepackage{graphicx} % Allows including images
\usepackage{booktabs} % Allows the use of \toprule, \midrule and \bottomrule in tables
\usepackage{mathpartir}

\newcommand\Ty{\mathtt{Ty}}
\newcommand\Ter{\mathtt{Ter}}
\newcommand\ttp{\mathtt{p}}
\newcommand\ttq{\mathtt{q}}
\newcommand\bfone{\mathbf{1}}
\newcommand\Set{\mathbf{Set}}
\newcommand\C{\mathbf{C}}
\newcommand\hatC{\widehat{\mathbf{C}}}
\newcommand\extop{{\scalebox{0.4}[1]{\_}}.{\scalebox{0.4}[1]{\_}}}
\DeclareMathOperator\obj{obj}

%----------------------------------------------------------------------------------------
%    TITLE PAGE
%----------------------------------------------------------------------------------------

\title{预层范畴上的 CwF 结构}
%\subtitle{Subtitle}

%\author{Pin-Yen Huang}

%\institute
%{
%    Department of Computer Science and Information Engineering \\
%    National Taiwan University % Your institution for the title page
%}
%\date{\today} % Date, can be changed to a custom date
\date{}

%----------------------------------------------------------------------------------------
%    PRESENTATION SLIDES
%----------------------------------------------------------------------------------------

\begin{document}

\begin{frame}
    % Print the title page as the first slide
    \titlepage
\end{frame}

\begin{frame}{概述}
  \begin{itemize}
    \item 任何预层范畴 $\hatC = \Set^{\C^{op}}$ 都可以诱导出一个 CwF 结构
    \item 对一个范畴 $\C$, 一般用 $I$, $J$, $K$ 表示对象, 用 $f$, $g$, $h$ 表示态射
  \end{itemize}
\end{frame}

\begin{frame}[shrink]{预层范畴到 CwF}
\begin{itemize}
  \item $\C$ 是一个小范畴, $\hatC=\Set^{\C^{op}}$ 是 $\C$ 上面的函子范畴即预层范畴
  \item 我们把上下文范畴取为 $\hatC$. 终对象就是把所有对象映为单元素集合的常量预层($\bfone(I)=\{\ast\}$, $\bfone(f) = id_{\{\ast\}}$. 取一个确定的单元素集合就好, 此处我们取$\{\ast\}$)
  \item 令 $\Delta$, $\Gamma$ 是 $\C$ 上的两个预层, $\gamma:\Delta\to\Gamma$ 是它们之间的一个自然变换, 那么我们有元素范畴之间的一个函子 $\int\gamma:\int\Delta\to\int\Gamma$. $(\int\gamma)(I,\delta) = (I,\gamma_{I}(\delta))$, $(\int\gamma)(f) = f$. 自然得到一个函子 $(\int\gamma)^{op}:(\int\Delta)^{op}\to(\int\Gamma)^{op}$. 函子 $\Ty:\hatC^{op}\to\Set$ 可以定义如下: $\Ty(\Gamma) \triangleq \obj(\Set^{(\int\Gamma)^{op}})$, 对 $\gamma:\Delta\to\Gamma$, $\Ty(\gamma):\Ty(\Gamma)\to\Ty(\Delta)$, $\Ty(\gamma)(A) \triangleq A \circ (\int\gamma)^{op}$. 我们也把$\Ty(\gamma)(A)$记为$A[\gamma]$. 那么我们有 $A[\gamma](I,\delta) = A(I,\gamma_I(\delta))$. 对类型 $A\in\Ty(\Gamma)$, 我们有一族集合$\{A(I,\rho) \mid I\in\C, \rho\in\Gamma(I)\}$, 使得对任意 $f : J \to I$, 我们有 $A(f) : A(I,\rho) \to A(J,\Gamma(f)(\rho))$, 满足 $A(id) = id$, $A(g\circ f) = A(f) \circ A(g)$. ($\int\Gamma$ 中的态射 $\bar{f} : (J,\Gamma(f)(\rho)) \to (I, \rho)$)
\end{itemize}
\end{frame}

\begin{frame}[shrink]{预层范畴到 CwF(续)}
\begin{itemize}
  \item 函子 $\Ter:(\int\Ty)^{op}\to\Set$ 定义如下: $\Ter(\Gamma\vdash A)\triangleq\hom(\bfone,A)$. 令 $a\in\Ter(\Gamma\vdash A)$, 那么 $a$ 是 $\bfone$ 和 $A$ 这两个 $\int\Gamma$ 上的预层之间的自然变换. 对任意$I\in\C$, $\rho\in\Gamma(I)$, 我们有 $a_{(I,\rho)} : \bfone(I,\rho) \to A(I,\rho)$, $a_{(I,\rho)} : \{\ast\} \to A(I,\rho)$, 我们把 $a_{(I,\rho)}(\ast)$ 简写为 $a(I,\rho)$. 这样也可以把 $a$ 看作是一族元素 $\{a(I,\rho)\in A(I,\rho) \mid I\in\C, \rho\in\Gamma(I)\}$, 使得对任意 $f:J\to I$, 我们有 $A(f)(a(I,\rho)) = a(J,\Gamma(f)(\rho))$. 对 $\int\Ty$ 中的一个态射 $\bar{\gamma}:(\Delta,A')\to(\Gamma,A)$, 我们有 $\gamma:\Delta \to \Gamma$ 是 $\hatC$ 中的一个态射, 并且 $A'=\Ty(\gamma)(A)$. $\Ter(\bar{\gamma}) : \Ter(\Gamma,A) \to \Ter(\Delta,A')$, $\Ter(\bar{\gamma})(a)(I,\delta)\triangleq a(I,\gamma_{I}(\delta))\quad\in\quad A(I,\gamma_{I}(\delta))=\Ty(\gamma)(A)(I,\delta)=A'(I,\delta)$
\end{itemize}
\end{frame}

\begin{frame}[shrink]{预层范畴到 CwF(续)}
\begin{itemize}
  \item 给定$\Gamma\in\hatC$和$A\in\Ty(\Gamma)$, 扩展上下文$\Gamma.A\in\hatC$定义如下
  \[ (\Gamma.A)(I) \triangleq \{ (\rho,a) \mid \rho \in \Gamma(I), a \in A(I,\rho) \} \]
  对态射 $f : J \to I$, $(\Gamma.A)(f) : (\Gamma.A)(I) \to (\Gamma.A)(J)$, $(\Gamma.A)(f) : \{(\rho,a) \mid \rho \in \Gamma(I), a \in A(I,\rho) \} \to \{(\rho,a) \mid \rho \in \Gamma(J), a \in A(J,\rho) \}$, $(\Gamma.A)(f)(\rho,a)\triangleq (\Gamma(f)(\rho),A(f)(a))$. $\bar{f} : (J,\Gamma(f)(\rho)) \to (I,\rho)$, $A(\bar{f}) : A(I,\rho) \to A(J,\Gamma(f)(\rho))$
  \item $\ttp : \Gamma.A \to \Gamma$, $\ttp_I : (\Gamma.A)(I) \to \Gamma(I)$, $\ttp_I : \{ (\rho,a) \mid \rho \in \Gamma(I), a \in A(I,\rho) \} \to \Gamma(I)$
  \[ \ttp_I(\rho,a) \triangleq \rho \]
  \item $\ttq \in \Ter(\Gamma.A \vdash A[\ttp])$, $\ttq$ 是 $\int\Gamma.A$ 上的两个预层 $\bfone$ 和 $A[\ttp]$ 之间的自然变换. $\ttq_{(I,(\rho,a))} : \bfone(I,(\rho,a)) \to (A[\ttp])(I,(\rho,a))$, $\ttq_{(I,(\rho,a))} : \{\ast\} \to A(I,\ttp_I(\rho,a))$, $\ttq_{(I,(\rho,a))} : \{\ast\} \to A(I,\rho)$, $\ttq(I,(\rho,a)) \triangleq \ttq_{(I,(\rho,a))}(\ast) \in A(I,\rho)$
  \[ \ttq(I,(\rho,a)) \triangleq a \]
\end{itemize}
\end{frame}


\begin{frame}[shrink]{预层范畴到 CwF(续)}
\begin{itemize}
  \item \phantom{ }\inferrule{\gamma:\Delta\to\Gamma \\ \Gamma\vdash A \\ \Delta\vdash a:A[\gamma]}{\langle\gamma,a\rangle:\Delta\to\Gamma.A}
  \[ \langle\gamma,a\rangle_I(\delta) \triangleq (\gamma_I(\delta),a(I,\delta)) \]
  \[ (\Gamma.A)(I) \triangleq \{ (\rho,a) \mid \rho \in \Gamma(I), a \in A(I,\rho) \} \]
  \[ a(I,\delta) \in A[\gamma](I,\delta) = A(I,\gamma_I(\delta)) \]
  \[ (\gamma_I(\delta),a(I,\delta)) \in (\Gamma.A)(I) \]
\end{itemize}
\end{frame}

\begin{frame}{额外说明}
  我们从预层范畴得到的这种形式的 CwF, 可以解释依赖和, 依赖积, 严格内蕴相等类型和宇宙
\end{frame}

\end{document}